\iffalse
\lstset{  % listings
    language=Python,
    basicstyle=\ttfamily\small,
    keywordstyle=\color{blue}\bfseries,
    stringstyle=\color{red},
    commentstyle=\color{green!60!black}\itshape,
    showstringspaces=false,
    frame=single,
    numbers=left,
    numberstyle=\tiny\color{gray},
    breaklines=true
}
\fi

\definecolor{yamlvalue}{HTML}{0550AE}
\definecolor{yamlkey}{HTML}{24292F}
\definecolor{yamlkeyword}{HTML}{CF222E}
\definecolor{yamlcomment}{HTML}{57606A}
\definecolor{yamlnumber}{HTML}{8C95A0}
\definecolor{yamlbackground}{HTML}{F6F8FA}
\definecolor{yamlline}{HTML}{D0D7DE}

\newcommand\YAMLkeystyle{\color{yamlkey}\bfseries}
\newcommand\YAMLvaluestyle{\color{yamlvalue}\mdseries}
\newcommand\YAMLcolonstyle{\color{yamlvalue}\mdseries}

\newcommand\ProcessThreeDashes{\textcolor{yamlcomment}{---}}

\lstdefinelanguage{yaml}{
    sensitive=true,
    morecomment=[l]{\#},
    morestring=[b]",
    morestring=[b]',
    moredelim=**[il][\YAMLcolonstyle{:}\YAMLvaluestyle]{:},
    literate={---}{{\ProcessThreeDashes}}3,
    keywords={true, false, null, yes, no, on, off},
    keywordstyle=\color{yamlkeyword}\bfseries
}

\lstdefinestyle{yamlstyle}{
    language=yaml,
    basicstyle=\small\ttfamily\YAMLkeystyle,
    columns=fullflexible,
    keepspaces=true,
    numbers=left,
    numberstyle=\tiny\ttfamily\color{yamlnumber},
    stepnumber=1,
    numbersep=12pt,
    backgroundcolor=\color{yamlbackground},
    frame=single,
    rulecolor=\color{yamlline},
    framerule=0.8pt,
    framesep=8pt,
    xleftmargin=24pt,
    xrightmargin=8pt,
    upquote=true,
    breaklines=true,
    breakatwhitespace=true,
    showstringspaces=false,
    upquote=true,
    tabsize=1,
    showlines=false,
    emptylines=0,
}

% ============================================================
% Palette Colori "Python Modern"
% ============================================================
\definecolor{pykeyword}{HTML}{CF222E}   % Rosso per keyword (def, class, import, return)
\definecolor{pybuiltin}{HTML}{953097}   % Viola per funzioni built-in (print, len, range)
\definecolor{pystring}{HTML}{0A3069}    % Blu scuro per le stringhe
\definecolor{pycomment}{HTML}{57606A}   % Grigio per i commenti
\definecolor{pynumber}{HTML}{005A9C}    % Blu brillante per i numeri
\definecolor{pybg}{HTML}{F6F8FA}        % Sfondo grigio chiarissimo
\definecolor{pyline}{HTML}{D0D7DE}      % Bordo sottile elegante
\definecolor{pyclass}{HTML}{8250DF}     % Viola per le classi
\definecolor{pymethod}{HTML}{0550AE}    % Blu per i metodi
\definecolor{pydecorator}{HTML}{B08800} % Giallo/arancio per i decoratori

% ============================================================
% Definizione dello Stile Python
% ============================================================
\lstdefinestyle{pystyle}{
    language=Python,
    basicstyle=\small\ttfamily\color{black},
    keywordstyle=\color{pykeyword}\bfseries,
    stringstyle=\color{pystring},
    commentstyle=\color{pycomment}\itshape,
    moredelim=**[is][\color{pynumber}]{£}{£},
    backgroundcolor=\color{pybg},
    frame=single,
    rulecolor=\color{pyline},
    framerule=0.8pt,
    framesep=8pt,
    xleftmargin=24pt,
    xrightmargin=8pt,
    emph=[1]{print,len,range,str,int,float,dict,list,set,enumerate,zip},
    emphstyle=[1]\color{pymethod}, % Evidenziazione avanzata per numeri, built-in e decoratori
    emph=[2]{BotnetInterface,Botnet},
    emphstyle=[2]\color{pymethod}\bfseries, % Evidenziazione avanzata per classi
    emph=[3]{connect,sniff},
    emphstyle=[3]\color{pymethod}\bfseries, % Evidenziazione avanzata per metodi
    % emph=[4]{@abstractmethod},
    % emphstyle=[4]\color{pydecorator}\bfseries, % Evidenziazione avanzata per decoratori
    numbers=left,
    numberstyle=\tiny\ttfamily\color{pycomment},
    numbersep=12pt,
    breaklines=true,
    breakatwhitespace=true,
    showstringspaces=false,
    upquote=true,
    tabsize=1
}
